\chapter{Introduction}
\label{chap:introduction}

Wildfires are becoming an increasingly significant source of economic, ecological, and public-health disruption in many regions of the world. In recent years, machine learning and deep learning have substantially improved the ability to predict wildfire occurrence, susceptibility, spread, and short-term fire danger. However, a major part of the literature still treats wildfire hazard prediction as an endpoint rather than as an input to broader consequence analysis. As a result, there remains a gap between wildfire prediction and wildfire-loss assessment, especially when the objective is to estimate impacts in monetary units or to connect fire-related conditions with broader ecological consequences.

This thesis is motivated by that gap. The central idea is that wildfire risk should be viewed not only as the probability of fire occurrence, but also as a function of what is exposed, how vulnerable it is, and what losses may follow. In this sense, modern predictive modeling should support not only hazard estimation but also consequence-oriented analysis. This perspective is increasingly emphasized in recent literature on wildfire expected annual loss, economy-wide disruption, smoke-related health burdens, and ecological vulnerability \citep{asif2026supranational,wang2021economicfootprint,law2025anthropogenic,arrogantefunes2024ecological}.

The present work focuses on county-level wildfire-loss modeling using publicly accessible datasets. As a first experimental direction, wildfire expected annual loss from the FEMA National Risk Index is used as a monetary target variable. This target is then modeled using county-level demographic, exposure, and resilience indicators, and later augmented with environmental predictors derived from publicly accessible geospatial datasets. In this way, the thesis develops a scalable baseline that connects wildfire-relevant environmental conditions with a consequence-oriented loss metric.

The main objectives of the thesis are:
\begin{enumerate}
    \item to review recent literature at the intersection of wildfire prediction, wildfire-loss assessment, and emerging AI methods;
    \item to construct a county-level wildfire-loss dataset in monetary units using accessible public data;
    \item to evaluate whether county-level wildfire expected annual loss can be predicted from exposure, vulnerability, and environmental variables;
    \item to assess how much environmental augmentation improves the predictive performance and geographic robustness of the resulting models.
\end{enumerate}

The novelty of this thesis is not limited to a particular study territory. Instead, the thesis proposes a general consequence-oriented wildfire-risk modeling workflow that can be adapted to other regions when comparable hazard, exposure, environmental, vulnerability, and loss-related data are available. The empirical implementation uses FEMA NRI wildfire Expected Annual Loss because it provides a standardized public monetary loss target, but the methodological logic is broader than this specific dataset.

Unlike much previous wildfire-machine-learning work that treats wildfire occurrence, susceptibility, fire danger, or burned area as the final output, this thesis uses wildfire prediction as a step toward loss-oriented risk assessment. The new contribution is to connect environmental wildfire drivers, exposure and value indicators, social vulnerability, community resilience, and monetary expected loss in one machine-learning pipeline.

The main contributions of the thesis are as follows:

\begin{enumerate}
    \item It provides a novel consequence-oriented framing of wildfire-risk modeling by shifting the focus from hazard prediction alone toward wildfire Expected Annual Loss estimation.
    
    \item It develops a new end-to-end machine-learning workflow that combines public risk indicators, exposure variables, vulnerability and resilience measures, topography, land cover, and climate predictors.
    
    \item Unlike models based only on socioeconomic or exposure variables, the thesis demonstrates that environmental augmentation substantially improves wildfire-loss prediction.
    
    \item It provides new insight into geographic transferability by comparing conventional random validation with state-wise grouped validation.
    
    \item It adds a novel exploratory ecological risk-screening layer that links predicted wildfire risk with natural vegetation exposure, without claiming to directly estimate ecological loss.
    
    \item It shows that the same general workflow can be extended beyond the current empirical territory to other regions, scales, and loss definitions, including event-level losses, asset-level exposure, ecological vulnerability, and post-fire recovery indicators.
\end{enumerate}

The experimental contribution of the thesis is a county-level wildfire Expected Annual Loss modeling pipeline built from FEMA National Risk Index data, followed by environmental augmentation using topography, land cover, and warm-season climate predictors. The results show that county-level wildfire loss can be modeled with meaningful predictive power and that environmental variables substantially improve both predictive performance and geographic transferability.

The implemented empirical workflow is intentionally county-level and national-scale, rather than event-level, pixel-level, or parcel-level. This choice enables reproducible modeling with accessible public data, while fine-scale wildfire tensor construction is retained as a future extension in Appendix~\ref{app:creek_fire_pipeline}.

Although the broader thesis theme includes both economic and ecological wildfire consequences, the implemented empirical pipeline focuses primarily on monetary wildfire Expected Annual Loss because FEMA NRI provides a standardized, publicly accessible county-level loss target. Ecological losses are not directly modeled as a separate damage target; however, the thesis adds an exploratory ecological risk-screening layer using natural vegetation exposure and predicted wildfire risk. This proxy-based analysis should be interpreted as ecological exposure/risk screening rather than direct ecological-loss estimation.

In addition to the main FEMA/NRI county-level pipeline, the thesis also briefly reports key results from a previous Vietnam case study on wildfire risk prediction in An Giang Province. That study used a compact deep neural network with multi-source environmental predictors, fire-growth feature engineering, threshold optimization, and SHAP-based interpretability. It is included as supporting previous work because it demonstrates the author's earlier experience with fine-scale wildfire prediction, while the main empirical contribution of this thesis remains the FEMA/NRI county-level wildfire Expected Annual Loss assessment pipeline.

Part of this work was successfully presented at the 4th International Conference on Distributed Computing and High Performance Computing (DCHPC 2026), held on 10--11 May 2026 in Tehran, Iran. The presented paper was entitled ``Mapping Forest Fire Risks Using Deep Learning: A Case Study in An Giang Province, Vietnam.'' The corresponding conference certificate is included in Appendix~\ref{app:conference_certificate}. The associated paper is expected to appear as a Web of Science/Scopus paper.

The remainder of the thesis is organized as follows. Chapter~\ref{chap:litreview} reviews the recent literature on wildfire risk, loss, ecological vulnerability, and emerging AI methods. Chapter~\ref{chap:methodology} describes the dataset construction process, predictor engineering, modeling workflow, and evaluation strategy. Chapter~\ref{chap:results} presents the experimental results and risk assessment. Chapter~\ref{chap:discussion} discusses the implications and limitations of the findings. Finally, Chapter~\ref{chap:conclusion} concludes the thesis and outlines directions for future work.